\chapter{선형사상, Kernel, Image}

\chapterintro{이 장에서는 벡터 자체의 계산에서 한 걸음 더 나아가 벡터공간 사이의 함수인 선형사상(linear map)을 중심으로 내용을 전개하겠습니다. 선형사상의 핵(kernel)과 상(image)을 통해 정보가 사라지는 방향과 전달되는 방향을 분리해 해석하고, 마지막에는 랭크-널리티(rank--nullity) 정리로 차원 관계를 완성하겠습니다.}
\chaptergoals{선형사상의 정의를 선형결합 보존 관점에서 정확히 다룰 수 있다.}{핵과 상을 계산하고 부분공간 구조 및 단사/전사 판정을 수행할 수 있다.}{합성, 역사상, 랭크-널리티 정리를 연결하여 유한차원 선형사상의 구조를 해석할 수 있다.}

\sectiontemplate{선형사상의 정의와 선형결합 보존}{선형사상은 덧셈과 스칼라곱을 동시에 보존하는 함수입니다. 이 절에서는 정의를 단순 암기가 아니라 계산 가능한 판정식으로 바꾸어 정리하겠습니다.}{\item 선형사상과 선형연산자(linear operator)를 구분할 수 있습니다.\item 선형결합 보존 성질을 이용해 선형성을 판정할 수 있습니다.\item 비선형 함수가 어떤 조건에서 실패하는지 반례로 설명할 수 있습니다.}{\item 벡터공간 공리\item 함수의 기본 성질과 합성}

\begin{definition}
체 $K$ 위 벡터공간 $V,W$ 사이의 함수 $T:V\to W$가 임의의 $u,v\in V$, $a\in K$에 대하여
\[
T(u+v)=T(u)+T(v),\qquad T(av)=aT(v)
\]
를 만족하면 $T$를 \emph{선형사상}(linear map)이라 한다.
\end{definition}

\begin{definition}
선형사상 $T:V\to V$처럼 정의역과 공역이 같은 경우를 \emph{선형연산자}(linear operator)라 한다.
\end{definition}

\begin{theorem}[선형결합 보존 판정]
함수 $T:V\to W$에 대하여 다음 두 조건은 동치이다.
\begin{enumerate}[label=(\roman*)]
\item $T$는 선형사상이다.
\item 임의의 자연수 $n$, 임의의 벡터 $v_1,\dots,v_n\in V$, 임의의 스칼라 $a_1,\dots,a_n\in K$에 대해
\[
T\!\left(\sum_{i=1}^n a_i v_i\right)=\sum_{i=1}^n a_i T(v_i)
\]
가 성립한다.
\end{enumerate}
\end{theorem}

\paragraph{증명 전략.}
정방향에서는 선형성의 두 공리를 이용해 유한합 공식을 귀납적으로 확장합니다. 역방향에서는 특별한 계수를 대입하여 덧셈 보존과 스칼라곱 보존을 각각 복원합니다.

\begin{proof}
(i)$\Rightarrow$(ii)를 보인다. $n=1$이면
\[
T(a_1v_1)=a_1T(v_1)
\]
이므로 성립한다. 이제 어떤 $n$에 대해 식이 성립한다고 가정하자. 그러면
\[
\begin{aligned}
T\!\left(\sum_{i=1}^{n+1}a_iv_i\right)
&=T\!\left(\left(\sum_{i=1}^{n}a_iv_i\right)+a_{n+1}v_{n+1}\right)\\
&=T\!\left(\sum_{i=1}^{n}a_iv_i\right)+T(a_{n+1}v_{n+1})\\
&=\sum_{i=1}^{n}a_iT(v_i)+a_{n+1}T(v_{n+1})\\
&=\sum_{i=1}^{n+1}a_iT(v_i).
\end{aligned}
\]
귀납법에 의해 모든 $n$에 대해 성립한다.

(ii)$\Rightarrow$(i)를 보인다. $n=2$, $(a_1,a_2)=(1,1)$, $(v_1,v_2)=(u,v)$를 대입하면
\[
T(u+v)=T(u)+T(v)
\]
를 얻는다. 또한 $n=1$, $a_1=a$, $v_1=v$를 대입하면
\[
T(av)=aT(v)
\]
를 얻는다. 따라서 $T$는 선형사상이다.
\end{proof}

\paragraph{의미 해석.}
이 정리는 선형성을 검사할 때 두 공리를 매번 따로 쓰지 않고, 하나의 선형결합 공식으로 한 번에 판정할 수 있게 해 줍니다. 이후의 계산은 이 공식을 기본 도구로 사용하면 크게 단순해집니다.

\begin{example}
$T:K^3\to K^2$를
\[
T(x,y,z)=(x-2y+z,\,3x+y)
\]
로 두면 $T$는 각 성분이 $(x,y,z)$의 선형결합이므로 선형사상이다. 반면
\[
S(x,y,z)=(x-2y+z+1,\,3x+y)
\]
는 $S(0,0,0)=(1,0)\ne(0,0)$이므로 선형사상이 아니다.
\end{example}

\subsection*{주의}
\begin{warning}
$T(0)=0$은 선형사상의 필요조건일 뿐 충분조건이 아니다. $T(v)=\|v\|u_0$ 같은 함수는 $T(0)=0$이지만 일반적으로 선형이 아니다.
\end{warning}
\begin{warning}
아핀 함수 $v\mapsto Av+b$에서 $b\ne 0$이면 선형성이 즉시 깨진다. 상수항 존재 여부를 먼저 확인해야 한다.
\end{warning}

\subsection*{자가진단퀴즈}
\begin{enumerate}[label=\arabic*.]
\item $T:\mathbb{R}^3\to\mathbb{R}^2$, $T(x,y,z)=(x+y,\,2y-z)$가 선형인지 정의로 판정하시오.
\item 선형사상 $T$에 대해 $T(-v)=-T(v)$, $T(u-v)=T(u)-T(v)$를 증명하시오.
\item $T(0)=0$이지만 선형이 아닌 함수의 구체적 예를 하나 구성하시오.
\end{enumerate}

\sectiontemplate{핵과 상: 부분공간 해석}{핵과 상은 선형사상의 구조를 가장 압축적으로 보여 주는 두 집합입니다. 핵은 사라지는 방향을, 상은 실제로 도달 가능한 출력의 범위를 나타냅니다.}{\item $\ker T$와 $\operatorname{im}T$를 정확히 정의할 수 있습니다.\item 핵과 상이 부분공간임을 엄밀히 증명할 수 있습니다.\item 단사성 판정을 $\ker T$로 수행할 수 있습니다.}{\item 부분공간 판정\item 연립일차방정식의 해집합}

\begin{definition}
선형사상 $T:V\to W$에 대하여
\[
\ker T=\{v\in V: T(v)=0\},\qquad
\operatorname{im}T=\{T(v): v\in V\}
\]
를 각각 $T$의 \emph{핵}(kernel), \emph{상}(image)이라 한다.
\end{definition}

\begin{theorem}
선형사상 $T:V\to W$에 대하여 다음이 성립한다.
\begin{enumerate}[label=(\roman*)]
\item $\ker T$는 $V$의 부분공간이다.
\item $\operatorname{im}T$는 $W$의 부분공간이다.
\item $T$가 단사일 필요충분조건은 $\ker T=\{0\}$이다.
\end{enumerate}
\end{theorem}

\paragraph{증명 전략.}
부분공간 판정은 ``영벡터 포함 + 선형결합 닫힘''으로 처리합니다. 단사 조건은 $T(u)=T(v)$를 $T(u-v)=0$으로 바꿔 핵 조건에 연결하면 깔끔하게 증명됩니다.

\begin{proof}
(i) 먼저 $0\in\ker T$임을 보인다. 선형성으로 $T(0)=T(0\cdot v)=0\cdot T(v)=0$이므로 성립한다.  
이제 $x_1,x_2\in\ker T$, $a,b\in K$를 잡으면
\[
T(ax_1+bx_2)=aT(x_1)+bT(x_2)=a\cdot 0+b\cdot 0=0
\]
이므로 $ax_1+bx_2\in\ker T$다. 따라서 $\ker T$는 부분공간이다.

(ii) $0=T(0)\in\operatorname{im}T$이므로 영벡터가 포함된다.  
$y_1,y_2\in\operatorname{im}T$이면 $y_1=T(v_1)$, $y_2=T(v_2)$인 $v_1,v_2\in V$가 존재한다. 임의의 $a,b\in K$에 대해
\[
ay_1+by_2=aT(v_1)+bT(v_2)=T(av_1+bv_2)\in\operatorname{im}T
\]
이므로 선형결합 닫힘이 성립한다. 따라서 $\operatorname{im}T$는 부분공간이다.

(iii) 먼저 $T$가 단사라고 가정하자. $v\in\ker T$이면 $T(v)=0=T(0)$이므로 단사성에 의해 $v=0$이다. 따라서 $\ker T=\{0\}$.  
역으로 $\ker T=\{0\}$라고 가정하자. $T(u)=T(v)$이면 선형성으로
\[
T(u-v)=T(u)-T(v)=0
\]
이므로 $u-v\in\ker T=\{0\}$, 따라서 $u=v$이다. 즉 $T$는 단사이다.
\end{proof}

\paragraph{의미 해석.}
핵과 상은 선형사상의 본질을 ``무엇을 없애는가''와 ``어디까지 도달하는가''로 분해해 보여 줍니다. 특히 단사성은 핵의 크기 하나로 즉시 판정되므로 계산과 이론이 자연스럽게 연결됩니다.

\begin{example}
$T:\mathbb{R}^3\to\mathbb{R}^2$를
\[
T(x,y,z)=(x+y,\;y+z)
\]
로 두자.  
$T(x,y,z)=(0,0)$에서 $x=-y$, $z=-y$이므로
\[
\ker T=\{(-t,t,-t): t\in\mathbb{R}\}
=\operatorname{span}\{(-1,1,-1)\}.
\]
또한 임의의 $(a,b)\in\mathbb{R}^2$에 대해 $(x,y,z)=(a,0,b)$를 택하면 $T(x,y,z)=(a,b)$이므로
\[
\operatorname{im}T=\mathbb{R}^2.
\]
\end{example}

\subsection*{주의}
\begin{warning}
$\ker T$는 정의역 $V$의 부분공간이고, $\operatorname{im}T$는 공역 $W$의 부분공간이다. 서로 다른 공간에 사는 집합을 같은 위치에서 비교하면 오류가 생긴다.
\end{warning}
\begin{warning}
$\operatorname{im}T=W$인지 확인하지 않고 전사라고 결론내리면 안 된다. 공역의 차원과 실제 상의 차원을 반드시 비교해야 한다.
\end{warning}

\subsection*{자가진단퀴즈}
\begin{enumerate}[label=\arabic*.]
\item $T:\mathbb{R}^4\to\mathbb{R}^2$, $T(x_1,x_2,x_3,x_4)=(x_1-x_2,\;x_3+x_4)$의 핵과 상을 구하시오.
\item ``$T$ 단사 $\Leftrightarrow \ker T=\{0\}$''를 정의만 이용해 다시 증명하시오.
\item 주어진 $w\in W$가 $\operatorname{im}T$에 속하는지 확인하는 절차를 선형시스템 관점으로 설명하시오.
\end{enumerate}

\sectiontemplate{합성, 선형동형, 역사상}{선형사상은 합성해도 선형이며, 가역인 경우 역사상도 선형으로 남습니다. 이 절에서는 선형동형(isomorphism)을 통해 ``공간은 다르지만 구조는 같다''는 관점을 엄밀히 정리하겠습니다.}{\item 합성사상의 선형성을 직접 증명할 수 있습니다.\item 선형동형사상의 정의와 의미를 설명할 수 있습니다.\item 전단사 선형사상의 역사상이 선형임을 증명할 수 있습니다.}{\item 함수 합성\item 전단사와 역함수}

\begin{definition}
벡터공간 $U,V,W$와 함수 $T:U\to V$, $S:V\to W$가 주어지면
\[
(S\circ T)(u)=S(T(u))\qquad (u\in U)
\]
로 정의되는 함수 $S\circ T:U\to W$를 \emph{합성}이라 한다.
\end{definition}

\begin{definition}
선형사상 $T:V\to W$가 전단사이면 $T$를 \emph{선형동형사상}(isomorphism)이라 하며, 이때 $V$와 $W$는 서로 \emph{동형}이라 한다.
\end{definition}

\begin{theorem}
벡터공간 $U,V,W$에 대하여 다음이 성립한다.
\begin{enumerate}[label=(\roman*)]
\item $T:U\to V$, $S:V\to W$가 선형사상이면 $S\circ T:U\to W$도 선형사상이다.
\item $T:U\to V$가 전단사 선형사상이면 역사상 $T^{-1}:V\to U$도 선형사상이다.
\end{enumerate}
\end{theorem}

\paragraph{증명 전략.}
합성의 경우에는 정의를 그대로 대입해 선형성 식을 단계적으로 전개하면 됩니다. 역사상의 선형성은 출력 벡터를 $T$의 원소로 되돌린 뒤, 다시 $T^{-1}$를 적용하는 방식으로 보입니다.

\begin{proof}
(i) $u_1,u_2\in U$, $a,b\in K$를 택하면
\[
\begin{aligned}
(S\circ T)(au_1+bu_2)
&=S(T(au_1+bu_2))\\
&=S(aT(u_1)+bT(u_2))\\
&=aS(T(u_1))+bS(T(u_2))\\
&=a(S\circ T)(u_1)+b(S\circ T)(u_2).
\end{aligned}
\]
따라서 $S\circ T$는 선형사상이다.

(ii) $T$가 전단사이므로 $T^{-1}$가 함수로 존재한다. $y_1,y_2\in V$, $a,b\in K$를 잡고
\[
x_1=T^{-1}(y_1),\qquad x_2=T^{-1}(y_2)
\]
라 두면 $y_1=T(x_1)$, $y_2=T(x_2)$이다. 선형성으로
\[
ay_1+by_2=aT(x_1)+bT(x_2)=T(ax_1+bx_2).
\]
양변에 $T^{-1}$를 적용하면
\[
T^{-1}(ay_1+by_2)=ax_1+bx_2
=aT^{-1}(y_1)+bT^{-1}(y_2).
\]
그러므로 $T^{-1}$는 선형사상이다.
\end{proof}

\paragraph{의미 해석.}
이 정리는 선형사상들의 세계가 합성과 역연산에 대해 안정적임을 보여 줍니다. 따라서 좌표변환, 기저변환, 동형분류를 모두 같은 대수적 틀 안에서 다룰 수 있습니다.

\begin{example}
$T:\mathbb{R}^2\to\mathbb{R}^2$를
\[
T(x,y)=(x+y,\;x-y)
\]
로 두면 $T$는 전단사 선형사상이고
\[
T^{-1}(u,v)=\left(\frac{u+v}{2},\frac{u-v}{2}\right)
\]
이다. 위 공식은 $T^{-1}$가 실제로 선형임을 직접 보여 준다.
\end{example}

\subsection*{주의}
\begin{warning}
합성의 순서는 중요하다. 일반적으로 $S\circ T$와 $T\circ S$는 동시에 정의되지도 않으며, 정의되더라도 서로 다를 수 있다.
\end{warning}
\begin{warning}
$T^{-1}$ 표기는 $T$가 전단사일 때만 사용할 수 있다. 단사 또는 전사 한쪽만 만족하면 역함수는 정의되지 않는다.
\end{warning}

\subsection*{자가진단퀴즈}
\begin{enumerate}[label=\arabic*.]
\item 선형사상 $T:\mathbb{R}^2\to\mathbb{R}^3$, $S:\mathbb{R}^3\to\mathbb{R}^2$를 임의로 하나씩 정하고 $S\circ T$를 계산하시오.
\item $S\circ T$가 단사이면 $T$가 단사임을 증명하시오.
\item $S\circ T$가 전사이면 $S$가 전사임을 증명하시오.
\end{enumerate}

\sectiontemplate{랭크-널리티 정리와 차원 공식}{유한차원 선형사상에서는 핵의 차원과 상의 차원이 정확히 보존식을 이룹니다. 이 절에서 랭크-널리티 정리를 완전증명하고, 차원 판정을 구조적으로 정리하겠습니다.}{\item nullity와 rank를 정의하고 계산할 수 있습니다.\item rank--nullity 정리를 완전하게 증명할 수 있습니다.\item 차원 정보로 단사/전사 여부를 판정할 수 있습니다.}{\item 기저와 기저 확장\item 부분공간의 차원}

\begin{definition}
$V$가 유한차원이고 $T:V\to W$가 선형사상일 때
\[
\operatorname{nullity}(T)=\dim\ker T,\qquad
\operatorname{rank}(T)=\dim\operatorname{im}T
\]
를 각각 $T$의 \emph{널리티}(nullity), \emph{랭크}(rank)라 한다.
\end{definition}

\begin{theorem}[Rank--Nullity]
$V$가 유한차원이고 $T:V\to W$가 선형사상이면
\[
\dim V=\operatorname{nullity}(T)+\operatorname{rank}(T)
\]
가 성립한다.
\end{theorem}

\paragraph{증명 전략.}
핵의 기저를 먼저 잡아 정의역의 기저로 확장합니다. 그다음 확장된 나머지 벡터들의 영상이 상의 기저가 됨을 보이면 차원 식이 즉시 따라옵니다.

\begin{proof}
$\ker T$의 기저를
\[
\{u_1,\dots,u_k\}
\]
라 하자. 기저 확장 정리에 의해 이를 $V$의 기저
\[
\{u_1,\dots,u_k,v_1,\dots,v_r\}
\]
로 확장할 수 있다. 따라서
\[
\dim V=k+r,\qquad k=\dim\ker T=\operatorname{nullity}(T).
\]
이제 $\{T(v_1),\dots,T(v_r)\}$가 $\operatorname{im}T$의 기저임을 보인다.

먼저 생성성을 보이자. 임의의 $y\in\operatorname{im}T$를 잡으면 어떤 $x\in V$가 존재하여 $y=T(x)$이다. $V$의 기저 전개로
\[
x=\sum_{i=1}^k a_i u_i+\sum_{j=1}^r b_j v_j
\]
라 쓰면
\[
\begin{aligned}
y=T(x)
&=\sum_{i=1}^k a_iT(u_i)+\sum_{j=1}^r b_jT(v_j)\\
&=\sum_{i=1}^k a_i\cdot 0+\sum_{j=1}^r b_jT(v_j)
=\sum_{j=1}^r b_jT(v_j).
\end{aligned}
\]
따라서 $\operatorname{im}T$는 $T(v_1),\dots,T(v_r)$로 생성된다.

다음으로 일차독립을 보이자. 만약
\[
\sum_{j=1}^r c_jT(v_j)=0
\]
이면
\[
T\!\left(\sum_{j=1}^r c_jv_j\right)=0
\]
이므로 $\sum_{j=1}^r c_jv_j\in\ker T$다. 따라서 어떤 $d_1,\dots,d_k\in K$가 존재하여
\[
\sum_{j=1}^r c_jv_j=\sum_{i=1}^k d_i u_i
\]
가 된다. 정리하면
\[
\sum_{i=1}^k (-d_i)u_i+\sum_{j=1}^r c_jv_j=0.
\]
그런데 $\{u_1,\dots,u_k,v_1,\dots,v_r\}$는 $V$의 기저이므로 일차독립이다. 따라서 모든 계수는 $0$이고 특히 $c_1=\cdots=c_r=0$이다. 즉 $T(v_1),\dots,T(v_r)$는 일차독립이다.

결론적으로
\[
\dim\operatorname{im}T=r=\operatorname{rank}(T).
\]
따라서
\[
\dim V=k+r=\operatorname{nullity}(T)+\operatorname{rank}(T).
\]
정리가 증명되었다.
\end{proof}

\paragraph{의미 해석.}
랭크-널리티 식은 정의역의 차원이 ``사라진 자유도(널리티)''와 ``전달된 자유도(랭크)''의 합으로 정확히 분해된다는 뜻입니다. 선형사상의 정보 흐름을 수치적으로 추적하는 핵심 보존식이라 할 수 있습니다.

\begin{example}
미분연산자 $D:P_3(K)\to P_2(K)$, $D(p)=p'$를 보자.  
$\ker D$는 상수다항식 전체이므로 $\dim\ker D=1$이다. 또한 임의의 $q(t)\in P_2(K)$에 대해
\[
p(t)=\int_0^t q(s)\,ds
\]
를 택하면 $p\in P_3(K)$이고 $p'=q$이므로 $\operatorname{im}D=P_2(K)$, 따라서 $\dim\operatorname{im}D=3$이다.  
즉
\[
\dim P_3=4=1+3=\operatorname{nullity}(D)+\operatorname{rank}(D).
\]
\end{example}

\subsection*{주의}
\begin{warning}
랭크-널리티 정리는 정의역이 유한차원일 때 진술된다. 무한차원에서는 차원 합 공식이 같은 형태로 즉시 성립하지 않는다.
\end{warning}
\begin{warning}
$\operatorname{rank}(T)=\dim W$이면 전사이지만, 단순히 $\operatorname{rank}(T)=\dim V$라는 정보만으로는 공역에 대한 전사성을 결론낼 수 없다.
\end{warning}

\subsection*{자가진단퀴즈}
\begin{enumerate}[label=\arabic*.]
\item $\dim V=7$, $\operatorname{rank}(T)=4$일 때 $\operatorname{nullity}(T)$를 구하시오.
\item $\dim V=\dim W=5$인 상황에서 $\operatorname{nullity}(T)=0$이면 왜 $T$가 동형사상인지 설명하시오.
\item $\operatorname{rank}(T)\le \dim W$를 상의 정의만 사용해 증명하시오.
\end{enumerate}

\chapterapplicationtemplate{오차 검출 부호에서 검사행렬 $H\in M_{m,n}(K)$는 선형사상 $T_H:K^n\to K^m$, $T_H(x)=Hx$를 정의합니다. 이때 실제 부호어 집합은 $\ker T_H$이고, 가능한 syndrome 전체는 $\operatorname{im}T_H$입니다. 따라서 $\dim\ker T_H=n-\operatorname{rank}(H)$라는 식을 이용하면, 중복검사 비트 수와 정보 비트 수의 균형을 차원 관점에서 직접 계산할 수 있습니다.}

\chaptersummarytemplate

\chapterexercisestemplate

\subsection*{연습문제 A (기초 확인)}
\begin{enumerate}[label=\arabic*.]
\item 다음 함수가 선형사상인지 판정하시오.  
(a) $T_1:\mathbb{R}^2\to\mathbb{R}^2$, $T_1(x,y)=(2x-y,\,3y)$  
(b) $T_2:\mathbb{R}^2\to\mathbb{R}$, $T_2(x,y)=x^2+y$  
(c) $T_3:P_2(\mathbb{R})\to P_2(\mathbb{R})$, $T_3(p)=p(1)\,t^2$
\item $T:\mathbb{R}^3\to\mathbb{R}^2$, $T(x,y,z)=(x+2y-z,\;y+z)$의 $\ker T$와 $\operatorname{im}T$를 구하시오.
\item 선형사상 $T$에 대해 $T(0)=0$, $T(-v)=-T(v)$를 증명하시오.
\item $T:\mathbb{R}^2\to\mathbb{R}^2$, $T(x,y)=(x+y,\;x-y)$의 역사상 공식을 구하고 선형성을 확인하시오.
\item $\dim V=6$이고 $\operatorname{rank}(T)=2$일 때 $\operatorname{nullity}(T)$를 구하시오.
\item $\dim V=5$, $\dim W=3$인 선형사상 $T:V\to W$에 대해 $T$가 단사일 수 없는 이유를 차원 공식으로 설명하시오.
\end{enumerate}

\subsection*{연습문제 B (개념 연결)}
\begin{enumerate}[label=\arabic*.]
\item 선형사상 $T:U\to V$, $S:V\to W$에 대해
\[
\ker(S\circ T)=\{u\in U: T(u)\in\ker S\}
\]
임을 증명하고, $\ker T\subseteq\ker(S\circ T)$를 유도하시오.
\item 매개변수 $a\in\mathbb{R}$에 대해 $T_a:\mathbb{R}^3\to\mathbb{R}^3$,
\[
T_a(x,y,z)=(x+ay,\;y+az,\;z)
\]
가 주어져 있다. $T_a$의 단사성/전사성/동형 여부를 $a$에 따라 판정하시오.
\item $D:P_4(\mathbb{R})\to P_3(\mathbb{R})$, $D(p)=p'$에 대해 rank와 nullity를 계산하고, 전사성과 단사성을 판정하시오.
\item 선형사상 $T:V\to W$에 대해 다음 두 조건의 동치를 증명하시오.  
(i) $T$는 단사이다.  
(ii) $V$의 임의의 유한 일차독립 집합의 영상은 $W$에서 일차독립이다.
\end{enumerate}

\subsection*{연습문제 C (증명/종합)}
\begin{enumerate}[label=\arabic*.]
\item 유한차원 벡터공간 $V$와 선형사상 $T:V\to W$, 부분공간 $U\le V$에 대해
\[
\dim T(U)=\dim U-\dim(U\cap\ker T)
\]
를 증명하시오.
\item 유한차원 벡터공간들 사이의 선형사상 $S:V\to W$, $T:W\to X$에 대해
\[
\operatorname{rank}(T\circ S)\le \min\{\operatorname{rank}(T),\operatorname{rank}(S)\}
\]
를 증명하시오.
\end{enumerate}